\section{K2 Thinking 的技术特点}

\subsection{自我验证能力}

模型在每一步 answer 后可以自我验证（self-verification），判断当前结果是否正确，决定是否需要继续推理或回溯。

\subsection{与 General Reward Model 的结合}

在 RL 训练中，可以为每步 think-answer 定义独立的 reward signal，而不仅仅是最终答案的 outcome reward。这使得训练信号更加密集和稳定。

\subsection{本章小结}

K2 Thinking 的技术核心在于通过 RL 训练模型的 interleaved reasoning 能力，结合细粒度 reward 和自我验证机制。

